\section{圆盘瞬子模空间的定向}
\label{sec:orientation-disc-instantons}

本节从\emph{开层}的切空间定向问题开始；该层由边界落在\termLagrangianSubmanifold{}
$L$ 上的光滑伪全纯圆盘组成。带拉格朗日边界的全纯映射模空间一般没有典范定向，
而闭全纯映射的模空间却有。定向的研究涉及带边界的\termFamilyIndexTheory{}。
事实表明，拉格朗日子流形的（相对）spin 结构与定向问题有关（FOOO09, Sil98）。\personPenkaGeorgieva{} (Geo13) 对一般开\personGromov{}--\personWitten{}理论中的定向作了优美阐述。

第~\ref{sec:coherent-orientations-floer-moduli}节将考察拉格朗日对 $(L_0,L_1)$
的光滑\personFloer{}模空间之定向问题，随后讨论紧化模空间上相容定向的存在性。对于后者，
我们将看到拉格朗日子流形的（相对）spin 结构起关键作用。

下面的命题说明，边界拉格朗日子丛的一个平凡化决定切空间的定向。

\begin{proposition}
\label{prop:trivial-boundary-orients-disc-index}
考虑 $(D^2,\partial D^2)$ 上的\termComplexBundlePair{} $(E,\lambda)$。
假设 $\lambda$ 平凡。则 $\lambda$ 的每个平凡化都典范地诱导
$\operatorname{Index}\overline\partial_{(E,\lambda)}$ 的一个定向。
\end{proposition}

\begin{proof}
按一般构造，模空间 $\widetilde{\mathcal M}(L;\beta;J)$ 带有\termDeterminantLineBundle{}
\[
 \det(D\overline\partial_J)\longrightarrow
 \widetilde{\mathcal M}(L;\beta;J),
\]
其在 $w$ 处的纤维为
\[
 \Lambda^{\max}\ker D_w\overline\partial_J
 \otimes\bigl(\Lambda^{\max}\operatorname{coker}D_w\overline\partial_J\bigr)^*.
\]
当 $\widetilde{\mathcal M}(L;\beta;J)$ 横截，即
$\operatorname{coker}D_w\overline\partial_J=\{0\}$ 时，有
\[
 \ker D_w\overline\partial_J\cong T_w\widetilde{\mathcal M}(L;\beta;J),
\]
故该\termDeterminantLineBundle{}上的定向给出模空间的定向。（一般而言，
$\det(D\overline\partial_J)$ 的定向按 (FOn99) 的意义给出
$\widetilde{\mathcal M}(L;\beta;J)$ 上\termKuranishiStructure{}的定向。）

先说明如何用拉格朗日子流形的几何，尤其是 spin 结构，逐纤维地赋予定向。
考虑带非空边界 $\partial\Sigma$ 的紧\personRiemann{}曲面 $\Sigma$ 上的复丛对
$(E,\lambda)$。以 $h$ 表示 $\partial\Sigma$ 的连通分支数，并以 $S_i$
表示第 $i$ 个连通分支，$i=1,\ldots,h$。

固定每个 $S_i$ 的一个环形邻域，对任意 $0<\epsilon<1$ 把它记为
$A_i((1-\epsilon,1])\subset D_i$。记子集
\[
 C_i=S^1(1-\epsilon)\subset A_i((1-2\epsilon,1]).
\]
沿 $\bigcup_iC_i$ 收缩 $\Sigma$，并令 $\Sigma/{\sim}$ 为把每个 $C_i$
等同为一点的商空间。所得商空间带有节点曲线结构，使商映射
\[
 \operatorname{cont}:\Sigma\longrightarrow\Sigma/{\sim}
 \cong\Sigma'\cup\left(\bigcup_iD_i\right)
\]
成为全纯映射，其中 $\Sigma'$ 是\emph{无}边紧\personRiemann{}曲面。以
$p_i\in\Sigma'$ 表示第 $i$ 个节点，即圆盘 $D_i$ 的中心附着之处。

任取一个 $S_i$ 并记为 $S$。利用给定同构
$E|_S\cong\lambda|_S\otimes\mathbb C$，$\lambda\to S$ 的\emph{每个平凡化}都诱导复平凡化
\begin{equation}
 \Phi:E|_{T_{2\epsilon}}\longrightarrow T_{2\epsilon}\times\mathbb C^n.
 \label{eq:boundary-trivialization-T-epsilon}
\end{equation}
\footnote{译者注：原书在此首次使用 $T_{2\epsilon}$，但本段此前只定义了
$A_i((1-\epsilon,1])$、$A_i((1-2\epsilon,1])$ 与 $C_i$，并未定义
$T_{2\epsilon}$。译文照录该符号。}
因此，$\Sigma$ 上的丛 $E$ 下降为节点曲线
\[
 \Sigma'\cup\left(\bigcup_iD_i\right)
\]
上的丛对 $(E',\lambda')$。

反过来，可解消节点并得到由
$\vec r=(r_1,\ldots,r_h)$、$r_i\in\mathbb R$ 参数化的一族\personRiemann{}曲面，
使丛对 $(E',\lambda')$ 本身形变为 $(\Sigma_r,\partial\Sigma_r)$ 上的丛对
$(E_r,\lambda_r)$，并有典范等同 $\lambda'\cong\lambda_r$。这是可能的，因为
$E'$ 在 $D_i$ 上\emph{已经平凡化}，并且给定了等同
\[
 E'_{\mathrm{ct},p_i}\cong E'_{\mathrm{end}_i,0}
 \cong\mathbb R^n\times\mathbb C,
 \qquad \lambda|_{\partial D_i}\cong\mathbb R^n,
\]
同时我们使用实的\emph{解消}参数 $r$。
\footnote{译者注：这里纤维记号 $\mathbb R^n\times\mathbb C$ 与复秩 $n$
的丛纤维类型并不相符；同页其余公式使用 $\mathbb C^n$。译文保留原式。}

该丛对的一个截面是二元组 $(\xi_{\mathrm{ct}},\xi_{\mathrm{end}})$，其中
$\xi_{\mathrm{end}}=(\xi_1,\ldots,\xi_h)$，每个 $\xi_i$ 是 $E'\to D_i$ 的截面，满足
$\xi_i(z)\in\lambda_{D_i,z}$ 及匹配条件
\begin{equation}
 \xi_{\mathrm{ct}}(p_i)=\xi_i(0).
 \label{eq:nodal-section-matching}
\end{equation}
换言之，$(E',\lambda')$ 的 $W^{1,p}$ 截面集为
\[
 W^{1,p}(E',\lambda'):=\Delta^{-1}(0),
\]
其中
\[
 \Delta:W^{1,p}(E'_{\mathrm{ct}})
 \oplus W^{1,p}(E'_{\mathrm{end}},\lambda')
 \longrightarrow\bigoplus_{i=1}^h\mathbb C_i^n
\]
定义为
\[
 \Delta(\xi_{\mathrm{ct}},\xi_{\mathrm{end}})
 =\bigl(\xi_{\mathrm{ct}}(p_1)-\xi_1(0),\ldots,
 \xi_{\mathrm{ct}}(p_h)-\xi_h(0)\bigr).
\]
这诱导正合列
\[
 0\to W^{1,p}(E',\lambda')
 \to W^{1,p}(E'_{\mathrm{ct}})\oplus W^{1,p}(E'_{\mathrm{end}},\lambda')
 \xrightarrow{\Delta}\bigoplus_{i=1}^h\mathbb C_i^n\to0.
\]
$\Sigma/{\sim}$ 上丛对 $(E',\lambda')$ 的相应指标，由算子
\[
 \overline\partial_{(E',\lambda')}:
 W^{1,p}(E',\lambda')\longrightarrow
 L^p(\Lambda^{0,1}(E'_{\mathrm{ct}}))
 \oplus L^p(\Lambda^{0,1}(E'_{\mathrm{end}}))
\]
的指标给出。

\begin{lemma}
\label{lem:boundary-trivialization-orients-nodal-index}
$\lambda\to\partial\Sigma$ 的每个平凡化都典范地决定
$\operatorname{index}\overline\partial_{(E',\lambda')}$ 的定向，因而也决定
$\det\overline\partial_{(E',\lambda')}$ 的定向。
\end{lemma}

\begin{proof}
先注意
\[
 \overline\partial_{(E'_{\mathrm{end}}|_{D_i},\lambda'|_{S_i})}:
 W^{1,p}(E'_{\mathrm{end}}|_{D_i},\lambda'|_{S_i})
 \longrightarrow L^p(\Lambda^{0,1}(E'_{\mathrm{end}}))
\]
是满射，而 $\lambda\to\partial\Sigma$ 的每个平凡化都诱导同构
\[
 \ker\overline\partial_{(E'_{\mathrm{end}}|_{D_i},\lambda'|_{S_i})}
 \cong\ker\overline\partial_{(D_i\times\mathbb C^n,S_i\times\mathbb R^n)}
 \cong\mathbb R^n.
\]
因此可取有限维\emph{复}子空间
\[
 \mathcal E\subset L^p(\Lambda^{0,1}(E')),
\]
使
\[
 (\mathcal E\oplus0)+\operatorname{Range}(\overline\partial_{(E',\lambda')})
 =L^p(\Lambda^{0,1}(E'_{\mathrm{ct}}))
 \oplus L^p(\Lambda^{0,1}(E'_{\mathrm{end}})).
\]
固定同构 $\Psi:\mathbb C^k\to\mathcal E\oplus0$，并考虑映射
\[
 \widehat{\overline\partial}_{(E',\lambda');\Psi}:
 \mathbb C^k\oplus W^{1,p}(E',\lambda')\longrightarrow
 \mathbb C^k\oplus
 \bigl(L^p(\Lambda^{0,1}(E'_{\mathrm{ct}}))
 \oplus L^p(\Lambda^{0,1}(E'_{\mathrm{end}}))\bigr),
\]
定义为
\[
 \widehat{\overline\partial}_{(E',\lambda');\Psi}(h,\xi)
 =\bigl(0,\Psi(h)+\overline\partial_{(E',\lambda')}(\xi)\bigr).
\]
\footnote{译者注：原书右端最后一项印作不含 $\xi$ 的
$\overline\partial_{(E',\lambda');\Psi}$，类型上不是向量；译文依该映射的定义域、
紧随其后的核方程及正合列补为
$\overline\partial_{(E',\lambda')}(\xi)$。}
有正合列
\[
 0\to\ker\overline\partial_{(E',\lambda')}
 \to\ker\widehat{\overline\partial}_{(E',\lambda');\Psi}
 \to\mathbb C^k\to\operatorname{coker}\overline\partial_{(E',\lambda')}\to0,
\]
其中中间三个映射依次为
\[
 d_1(\xi)=(0,\xi),\qquad d_2(h,\xi)=h,
 \qquad d_3(h)=\Psi(h)\bmod\operatorname{Range}\overline\partial_{(E',\lambda')}.
\]
还要注意，$\widehat{\overline\partial}_{(E',\lambda');\Psi}$ 通过\personFredholm{}算子
同伦于 $\operatorname{id}_{\mathbb C^k}\oplus\overline\partial_{(E',\lambda')}$。
故有
\begin{equation}
 \operatorname{Index}\overline\partial_{(E',\lambda')}
 \cong\operatorname{Index}\widehat{\overline\partial}_{(E',\lambda');\Psi}
 \cong\ker\widehat{\overline\partial}_{(E',\lambda');\Psi}-\mathbb C^k.
 \label{eq:stabilized-index-isomorphism}
\end{equation}
注意 $\ker\widehat{\overline\partial}_{(E',\lambda');\Psi}$ 是满足
\[
 0=\Psi(h)+\overline\partial_{(E',\lambda')}(\xi)
\]
的二元组 $(h,\xi)$ 之集，即
\[
 \ker\widehat{\overline\partial}_{(E',\lambda');\Psi}
 \cong\overline\partial_{(E',\lambda')}^{-1}(0\oplus\mathcal E).
\]
由于 $\mathcal E$ 的选取使取值映射
\[
 \operatorname{ev}_p:
 \overline\partial_{(E'|_{\mathbb CP^1})}^{-1}(\mathcal E)
 \longrightarrow E'_p\cong\mathbb C^n
\]
满射，$\overline\partial_{(E',\lambda')}^{-1}(0\oplus\mathcal E)$ 恰为满射同态
\[
\begin{aligned}
 \Delta:
 &\overline\partial_{(E'|_{\mathbb CP^1})}^{-1}(\mathcal E)
 \times\ker\overline\partial_{(E'|_D,\lambda'|_S)}\longrightarrow\mathbb C^n,\\
 &(\xi_{\mathbb CP^1},\xi_D)\longmapsto
 \xi_{\mathbb CP^1}(p)-\xi_D(0)
\end{aligned}
\]
的核。回忆，$\overline\partial_{(E'|_{\mathbb CP^1})}^{-1}(\mathcal E)$
和 $\mathbb C^n$ 是复向量空间，故具有典范定向。

这就证明：若给定每个 $\lambda$ 的平凡化，它将典范地决定
$\ker\overline\partial_{(E'|_D,\lambda')}$ 的定向，从而决定
$\overline\partial_{(E',\lambda')}^{-1}(0\oplus\mathcal E)$ 的定向；再由同构
\eqref{eq:stabilized-index-isomorphism}，它决定
$\operatorname{Index}\overline\partial_{(E',\lambda')}$ 的定向。
\end{proof}

接下来通过粘合把该定向转移到 $(E,\lambda)$。解消节点得到由 $r\in\mathbb R_+$
参数化的一族\personRiemann{}曲面，使丛对 $(E',\lambda')$ 形变为
$(\Sigma_r,\partial\Sigma_r)$ 上的丛对 $(E_r,\lambda_r)$，且有典范等同
$\lambda'\cong\lambda_r$。

更精确地，选取 $D^2$ 中 $0$ 和 $\mathbb CP^1$ 中 $p$ 的小坐标邻域
$D_\delta(0)$ 与 $D_{2\delta}(p)$。这里 $D_\delta(0)$ 是以 $0$ 为心、半径
$\delta$ 的圆盘。对正实数 $r>0$，在 $O$ 与 $S$ 附近粘合 $D^2$ 和
$\mathbb CP^1$：当 $zw=1/r$ 时，等同
$z\in D_{2\delta}^2(O)$ 与 $w\in D_{2\delta}^2(S)$。把所得带边\personRiemann{}曲面记作
$\Sigma_r$，它双全纯于单位圆盘。可以等同 $\Sigma_1=D^2$ 及
$\Sigma_\infty=\mathbb CP^1\cup D$。
\footnote{译者注：本段开头的节点记作 $0,p$，粘合公式却改用未在本段定义的
$O,S$；译文照录这组符号。}

由 $E'$ 还得到 $\Sigma_r$ 上的向量丛，记为 $E_r$。$\partial D$ 上的全实子丛
$\lambda'$ 诱导 $\partial\Sigma_r$ 上的全实子丛 $\lambda_r$。由于 $\mathcal E$
由支集远离节点的元素组成，它典范地嵌入
$L^p(\Lambda^{0,1}\Sigma_r\otimes E_r)$；故对任意充分大的 $r$，映射 $\Psi$
可典范地视为
\[
 \Psi:\mathbb C^k\longrightarrow L^p(\Lambda^{0,1}(\Sigma_r)\otimes E_r).
\]
于是由 $\widehat{\overline\partial}_{(E',\lambda');\Psi}$ 的满射性及粘合
(SS88)，有
\[
 \operatorname{Range}\widehat{\overline\partial}_{(E_r,\lambda_r);\Psi}
 +\mathcal E=L^p(\Lambda^{0,1}(\Sigma_r)\otimes E_r)
\]
及典范同构
\[
 \operatorname{Index}\widehat{\overline\partial}_{(E',\lambda');\Psi}
 \cong\operatorname{Index}\widehat{\overline\partial}_{(E_r,\lambda_r);\Psi}.
\]
因此，当 $r\to\infty$ 时
$\operatorname{Index}\widehat{\overline\partial}_{(E',\lambda');\Psi}$ 的定向，
典范地给出充分大 $r>0$ 时
$\operatorname{Index}\widehat{\overline\partial}_{(E_r,\lambda_r);\Psi}$ 的定向；
再由同构~\eqref{eq:stabilized-index-isomorphism}的第一部分，得到
$\operatorname{Index}\overline\partial_{(E_r,\lambda_r)}$ 的定向。随后把
$(E_r,\lambda_r)$ 连续形变为原给定丛对 $(E_1,\lambda_1)=(E,\lambda)$，
便在后者上诱导一个定向。命题得证。
\end{proof}

当模空间为零维、正则且可定向时，其中每个元素的定向之集同构于
$\mathbb Z_2$。模空间定向的一种选择，便给出该定向空间与 $\mathbb Z_2$
之间的同构。在维数 $4$ 时，可以用给定拉格朗日边界与伪全纯圆盘的几何相交，
显式描述模空间中元素的这一符号。（见 (Cho08)。）

\begin{example}[单点开 GW 不变量]
\label{ex:one-point-open-gw-invariant}
设 $L\subset M$ 是有向拉格朗日子流形。假设 $TL$ 平凡（例如，$L$ 为 spin），
并给定一个平凡化。
\footnote{译者注：原书括号中的 ``例如，$L$ 为 spin'' 并不成立于一般情形：
spin 性并不蕴含切丛平凡。译文保留原断言并在此指出。}
命题~\ref{prop:trivial-boundary-orients-disc-index} 给出模空间
\[
 \mathcal M_1(L,\beta)
 \cong\widetilde{\mathcal M}(L;\beta)\times\partial D^2/PSL(2,\mathbb R)
\]
的定向；这正是 $\widetilde{\mathcal M}(L;\beta)$ 与 $\partial D^2$ 的
纤维积定向，再取商。更具体地说，通过正合列
\begin{equation}
 0\to\mathfrak{psl}(2,\mathbb R)\cdot\{(w,z)\}
 \to T_w\widetilde{\mathcal M}(L;\beta)\oplus T_z\partial D^2
 \to T_{[w,z]}\mathcal M_1(L,\beta)\to0
 \label{eq:one-point-disc-orientation-sequence}
\end{equation}
定义 $\mathcal M_1(L,\beta)$ 的定向。

以 $\widehat a$ 表示由\personLie{}代数元素
$a\in\mathfrak{psl}(2,\mathbb R)$ 的线性化作用生成的 $D^2$ 上向量场。
回忆，$a\in\mathfrak{psl}(2,\mathbb R)$ 在 $(w,z)$ 上的线性化作用为
\begin{equation}
 a\longmapsto\bigl(-dw(\widehat a(z)),\widehat a(z)\bigr).
 \label{eq:linearized-psl-action-disc}
\end{equation}
令
\[
 \mathfrak{psl}(2,\mathbb R)=\mathfrak h\oplus\mathfrak p
 \cong\mathbb R\oplus\mathbb R^2
\]
为\personCartan{}分解，其中 $\mathfrak h$ 是对角矩阵的\personLie{}子代数。固定
$\mathfrak h$ 和 $\mathfrak p$ 上的定向，并在
$\mathfrak{psl}(2,\mathbb R)$ 上取直和定向。

作用~\eqref{eq:linearized-psl-action-disc}还诱导另一分裂
\[
 \mathfrak{psl}(2,\mathbb R)=\mathfrak{psl}(2,\mathbb R)_z\oplus\mathbb R,
\]
其中 $\mathfrak{psl}(2,\mathbb R)_z$ 是点 $z\in\partial D^2$ 的等变\personLie{}子代数，
它总包含 $\mathfrak h$。在上半空间模型中，该分裂中的 $\mathbb R$-作用对应于
$\partial/\partial x$ 生成的平移，并在 $\partial D^2$ 上诱导逆时针旋转。
以 $\partial/\partial x$ 固定 $\mathbb R$ 的定向；于是
$\mathfrak{psl}(2,\mathbb R)_z$ 带有该分裂诱导的典范定向。正合列
\eqref{eq:one-point-disc-orientation-sequence}中的映射
\[
 \mathfrak{psl}(2,\mathbb R)\cdot\{(w,z)\}
 \longrightarrow T_w\widetilde{\mathcal M}(L;\beta)\oplus T_z\partial D^2
\]
把这个 $\mathbb R$-因子映到 $T_z\partial D^2$。这些定向与
\eqref{eq:one-point-disc-orientation-sequence}一起，在每个
$[w,z]\in\mathcal M_1(L,\beta)$ 处诱导 $T_{[w,z]}\mathcal M_1(L,\beta)$ 的定向。

现假设 $\dim\mathcal M_1(L,\beta)=n$，并考虑取值映射
$\operatorname{ev}_1:\mathcal M_1(L,\beta)\to L$。由维数限制，对通有的
$x\in L$，可给 $\operatorname{ev}_1^{-1}(x)$ 中每个元素 $[w]$ 赋予符号，
记作 $\sigma([w])$。假设紧化模空间 $\overline{\mathcal M}_1(L,\beta)$ 横截，
即可把单点开\personGromov{}--\personWitten{}不变量定义为\footnote{译者注：原书先把符号
$\sigma([w])$ 赋给 $[w]\in\operatorname{ev}_1^{-1}(x)$，最后的求和指标却改写为
$[w]\in\mathcal M(L,\beta)$；后者并非前述零维取值纤维。译文照录求和指标。}
\[
 \sum_{[w]\in\mathcal M(L,\beta)}\sigma([w]).
\]
\end{example}
